% ============================================================
% Thesis Chapter: The UPF Digital Twin and its Construction
% ============================================================
% Compile with:  pdflatex chapter_digital_twin.tex  (run twice)
% Figures live in  reports/figures/  (PNGs produced by make_figures.py).
% Diagrams marked [DRAW.IO PLACEHOLDER] are to be drawn by hand; the box
% text is the brief for the drawing.
% ============================================================

\documentclass[12pt, a4paper]{report}

% -- Packages ------------------------------------------------------------
\usepackage[T1]{fontenc}
\usepackage[utf8]{inputenc}
\usepackage{lmodern}
\usepackage{amsmath, amssymb}
\usepackage{booktabs}
\usepackage{graphicx}
\usepackage{subcaption}
\usepackage{xcolor}
\usepackage[margin=2.5cm]{geometry}
\usepackage[hidelinks]{hyperref}
\usepackage{setspace}
\onehalfspacing

% -- Convenience macros --------------------------------------------------
\newcommand{\dpdk}{\textsc{sd-core dpdk}}
\newcommand{\usr}{\textsc{usr-upf}}
\newcommand{\upf}{\textsc{upf}}
\newcommand{\code}[1]{\texttt{#1}}
\newcommand{\figpath}[1]{figures/#1}

% Draw.io placeholder box: a framed brief describing a diagram to be drawn.
\newcommand{\drawioplaceholder}[1]{%
  \fbox{\parbox{0.94\linewidth}{\centering\vspace{2pt}
    \textbf{[\,DRAW.IO PLACEHOLDER\,]}\\[4pt]
    \footnotesize\raggedright #1\vspace{2pt}}}}

\begin{document}
% ============================================================

\chapter{The UPF Digital Twin}
\label{chap:digital_twin}

% -- Section 1: Introduction --------------------------------------------
\section{Introduction}
\label{sec:dt_intro}

Chapter~\ref{chap:upf_profiling} established, from controlled laboratory
measurements, \emph{how} two open-source \upf{} implementations behave under
load: \dpdk{} draws a near-constant 0.82\,W regardless of traffic, whereas
\usr{} scales its power proportionally to load but saturates---losing packets
and inflating delay---above roughly 0.5\,Gbps. Those measurements are the raw
material; on their own they do not yet constitute a tool a network controller
can query.

This chapter describes the construction of a \emph{digital twin}: a
measurement-grounded, controller-agnostic software component that answers a
single well-posed question,

\begin{quote}
  \emph{Given an offered traffic load and a candidate \upf{} configuration,
  what would happen---in terms of power, throughput, CPU, packet loss, delay,
  and QoS safety?}
\end{quote}

\noindent
It is important to be precise about what the twin \emph{is} and \emph{is not},
because the surrounding system is deliberately decomposed into independent
parts:

\begin{description}
  \item[The twin is not a traffic forecaster.] Where the offered load
    \emph{comes from}---a spatio-temporal forecast over base-station
    clusters---is the subject of its own chapter
    (Chapter~\ref{chap:forecaster}). The twin consumes a load value; it does
    not predict one.
  \item[The twin is not a controller.] \emph{Which} configuration to choose
    is a policy decision. The twin only evaluates the consequences of a
    decision that has already been made. Controllers---both rule-based
    baselines and learned (reinforcement-learning) agents---are developed and
    compared in Chapter~\ref{chap:controllers}. To keep the present chapter
    self-contained while avoiding any pre-emption of that material, every
    illustration here is driven by a single, transparent, \emph{non-learning}
    controller (a hysteresis rule, Section~\ref{sec:dt_eval}); the learned
    controllers and an informed oracle bound are reserved for
    Chapter~\ref{chap:controllers}.
\end{description}

This separation of concerns---forecaster, twin, controller---mirrors the
repository structure of the implementation, in which the twin is a
\code{pip}-installable library that the controller project depends on but
never modifies.

\paragraph{Place in the thesis.} Within this thesis's theme of \emph{augmented
observability} for the 5/6G data plane, the twin extends the observability
envelope along its \emph{counterfactual} axis, and this is the deepest extension
of the three. The profiling campaign made the present observable and the
forecaster made the near future observable, but both concern the state the system
is actually in. The twin makes observable the states it is \emph{not} in: the
power, QoS, and safety that \emph{would} result under the realisation that is not
currently running. Runtime control between two \upf{} realisations is impossible
to reason about without this---choosing requires comparing an outcome that is
being observed against one that never occurred---and it is exactly the signal the
controllers in Chapter~\ref{chap:controllers} consume when they weigh a switch.

The remainder of the chapter proceeds as follows.
Section~\ref{sec:dt_from_profiling} bridges the profiling results to a
predictive surrogate and justifies which trained variants the twin actually
uses. Section~\ref{sec:dt_arch} presents the two evaluation modes (stateless
physics and stateful rollout) and the data contract.
Section~\ref{sec:dt_inputs} describes load calibration and the
NetMob-compatible input adapter. Section~\ref{sec:dt_kpis} defines the
predicted KPIs and the safety/efficiency semantics.
Section~\ref{sec:dt_switching} formalises the switching-cost model.
Section~\ref{sec:dt_thresholds} shows how operating thresholds are
\emph{derived} from the surrogate rather than hardcoded.
Section~\ref{sec:dt_eval} defines the evaluation methodology and metrics and
demonstrates the assembled twin under the hysteresis controller.
Section~\ref{sec:dt_summary} summarises.

% -- Section 2: From profiling to a predictive surrogate ----------------
\section{From Profiling to a Predictive Surrogate}
\label{sec:dt_from_profiling}

The twin's predictive core is the two-layer surrogate introduced in
Section~\ref{sec:twin}. For completeness we restate its structure and then
focus on the design decisions specific to embedding it in an online evaluator.

\subsection{Two-layer structure}

For each \upf{} variant the surrogate is organised as a stack that mirrors the
causal pathway \emph{load $\to$ performance $\to$ power}:

\begin{description}
  \item[Layer~1 (performance).] Four independent single-target regressors map
    the offered load to delivered throughput (Gbps), CPU utilisation (\%),
    downlink packet loss (packets per interval), and downlink one-way delay
    ($\mu$s).
  \item[Layer~2 (power).] A single regressor maps the offered load
    \emph{plus} the four Layer~1 outputs to \code{power\_watts}.
\end{description}

Because each Layer~1 slot is selected independently by cross-validated model
search, the chosen algorithm differs by target: the near-linear
throughput--load relationship is best captured by Ridge regression, whereas
the strongly non-linear delay and loss responses near saturation are best
captured by tree ensembles (Random Forest or Gradient Boosting). This
heterogeneity is a feature, not an inconsistency: forcing a single algorithm
across all slots would underfit the non-linear targets. One practical
consequence, visible later in Figure~\ref{fig:dt_surrogate}, is that
tree-based slots produce \emph{piecewise-constant} predictions and do not
extrapolate beyond the measured envelope---a property the twin must respect
when it is queried.

\subsection{Which variants the twin uses}
\label{sec:dt_variants}

The profiling campaign trained three variants: \code{dpdk},
\code{usr\_full} (all USR samples, including the saturated regime), and
\code{usr\_safe} (USR samples restricted to the unsaturated regime below the
saturation knee). The twin exposes exactly two \emph{actions} to a
controller---\code{DPDK} and \code{USR}---and maps them to trained variants as
follows:
\[
  \texttt{DPDK} \longmapsto \texttt{dpdk}, \qquad
  \texttt{USR}  \longmapsto \texttt{usr\_full}.
\]

\noindent
The choice of \code{usr\_full} over \code{usr\_safe} to represent \code{USR}
is deliberate and central to the twin's purpose. A digital twin whose
remit explicitly includes \emph{packet loss, delay, and QoS safety} must be
able to represent what happens when \usr{} is pushed past its safe operating
point. \code{usr\_safe} was trained only on loads where \usr{} never loses
packets; queried above the knee it would extrapolate near-zero loss and low
delay, silently hiding the very failure mode the twin exists to detect.
\code{usr\_full}, having observed the saturation transition, predicts the
onset of loss and delay faithfully (the corresponding Layer~1 targets reach
$R^2 = 0.998$ and $0.968$ respectively, against $0.86$ and $0.52$ for
\code{usr\_safe}). The cost is a marginally noisier power model
($R^2 = 0.988$ vs.\ $0.999$), which is an acceptable price for correct
safety behaviour. As Section~\ref{sec:dt_thresholds} shows, the twin's
ability to \emph{derive} a QoS threshold depends entirely on this property:
with \code{usr\_safe} the threshold-derivation sweep would never observe an
unsafe step and the derived limit would be meaningless. The \code{usr\_safe}
models are therefore retained only as a low-load power reference and are not
wired into the twin's action map.

\subsection{Full and lite feature sets}

Each variant was trained in two feature regimes. The \emph{full} model accepts
all seven offered-load features (DL/UL kbits/s, TX packet counts, average
packet size, and the L2/L3 overhead ratio) and targets laboratory or
high-fidelity simulation. The \emph{lite} model accepts only the DL/UL load
volume and is therefore drivable from datasets---such as NetMob~2023---that
expose only aggregate per-cell load with no packet-level detail. The twin
loads the \emph{lite} variant exclusively, because its online use case is
driven by cluster-level load forecasts that carry no per-packet observability.
The accuracy sacrificed for this interoperability is small: for the power
model the lite--full gap is $\Delta R^2 \approx 0.024$ for \code{usr\_full}.

% -- Section 3: Architecture --------------------------------------------
\section{Twin Architecture}
\label{sec:dt_arch}

The twin offers two evaluation modes over the same surrogate core, reflecting
the two questions a consumer may ask (Figure~\ref{fig:dt_arch}).

\begin{description}
  \item[Stateless physics.] \code{evaluate\_action(action, load)} returns the
    predicted KPIs for a single (action, load) point. It is a pure function:
    no memory, no side effects. This is the mode used by parameter sweeps, the
    threshold derivation of Section~\ref{sec:dt_thresholds}, and interactive
    inspection.
  \item[Stateful rollout.] A \emph{session} wraps the stateless core with the
    previous action and applies the switching-cost accounting of
    Section~\ref{sec:dt_switching}. \code{session.step(requested\_action,
    load)} returns the realised KPIs for one time step, including any energy
    paid for activation. This is the mode used to drive a sequential
    simulation (and, in Chapter~\ref{chap:controllers}, the \code{step()} of
    a reinforcement-learning environment).
\end{description}

A consumer constructs a twin from a data directory that follows a fixed
contract: a \code{profiling\_twin/} subtree containing the surrogate
\code{models/} (with a \code{manifest.json} registry and \code{layer1/},
\code{layer2/} model files) and a \code{switching\_costs.yaml}. The twin is
otherwise configured by a scenario dictionary supplying the QoS budget, the
switching-accounting policy, and the time-step length. Crucially, the twin
holds no reference to any controller type; controllers depend on the twin, not
the reverse.

\begin{figure}[htbp]
  \centering
  \drawioplaceholder{%
    \textbf{Twin architecture / data flow.} A left-to-right block diagram.
    \emph{Inputs (left):} normalised cluster load $\hat d \in [0,1]$ from the
    forecaster (Chapter~\ref{chap:forecaster}); a requested action
    $\in\{\texttt{DPDK},\texttt{USR}\}$ from the controller
    (Chapter~\ref{chap:controllers}).
    \emph{Calibration block:} $\hat d \to$ load (Gbps) via the global
    $\alpha$; load (Gbps) $\to$ lite features (kbits/s).
    \emph{Surrogate core (centre):} two stacked boxes per action---``Layer~1:
    4 performance models'' feeding ``Layer~2: power model''---selected from the
    model registry (\texttt{dpdk} or \texttt{usr\_full}).
    \emph{Safety/efficiency block:} compares predicted delay/loss against the
    QoS budget ($\Rightarrow$ \texttt{is\_safe}) and predicted power against
    the alternative UPF ($\Rightarrow$ \texttt{is\_efficient}).
    \emph{Switching block (only in stateful mode):} previous action +
    activation duration + accounting rule $\Rightarrow$ switching energy.
    \emph{Outputs (right):} a \texttt{UPFResult} (power, delay, throughput,
    loss, \texttt{is\_safe}, \texttt{is\_efficient}) and, in stateful mode,
    the switching energy for the step. Draw the two modes as two horizontal
    lanes sharing the surrogate core; mark the stateless lane as ``pure
    function'' and the stateful lane as carrying state (previous action).}
  \caption{Architecture of the digital twin. The surrogate core is shared by
    a stateless physics mode and a stateful rollout mode; the latter adds the
    switching-cost model. The twin is controller-agnostic: forecaster and
    controller are external (Chapters~\ref{chap:forecaster}
    and~\ref{chap:controllers}).}
  \label{fig:dt_arch}
\end{figure}

% -- Section 4: Inputs and calibration ----------------------------------
\section{Inputs and Load Calibration}
\label{sec:dt_inputs}

\subsection{Normalised load to physical load}

The forecasting stage emits per-cluster load normalised to $[0,1]$. A single
global calibration constant $\alpha$ converts this to a physical offered load
in Gbps:
\begin{equation}
  \lambda_{\mathrm{Gbps}} = \alpha \cdot \hat d, \qquad
  \alpha = 1.0~\text{Gbps per normalised unit}.
  \label{eq:dt_alpha}
\end{equation}
The value of $\alpha$ is a scenario parameter; it sets the physical scale at
which the (otherwise dimensionless) forecast is interpreted, and hence where a
given normalised load falls relative to the \usr{} saturation knee.

It must be stated plainly that $\alpha$ is a \emph{declared modelling
assumption, not a recovered calibration}. The NetMob traces are dimensionless
by construction, for privacy reasons, and the forecaster sums per-base-station
normalised demand across each cluster; there is consequently no physical
scaler that could be looked up or fitted. Setting $\alpha = 1$ interprets one
unit of summed normalised demand as 1\,Gbps of offered \upf{} load, which
places the $K = 10$ clusters at mean loads of $0.11$--$1.66$\,Gbps with peaks
to $5.8$\,Gbps---a realistic per-site MEC \upf{} operating range. The same
value is used by the controller work of Chapter~\ref{chap:controllers}, so
that both chapters describe one physical regime. Section~\ref{sec:dt_results}
returns to the consequences of this choice, which are substantial: because the
profiled \usr{} data path saturates one to two orders of magnitude below these
loads, $\alpha$ largely determines how much room for energy-aware switching
exists at all.

\subsection{Lite input adapter}

The lite surrogate consumes load volume in kbits/s. The twin converts the
calibrated load directly,
\begin{equation}
  B = \lambda_{\mathrm{Gbps}} \times 10^{6}~[\text{kbits/s}],
  \label{eq:dt_kbits}
\end{equation}
and supplies $B$ as both the downlink and uplink lite features, matching the
NetMob compatibility adapter of Section~\ref{sec:netmob}. This is the only
transformation between a controller-facing load value and the model input;
all remaining (full-model) features are absent by construction in the lite
regime and are never required.

% -- Section 5: Predicted KPIs and safety semantics ---------------------
\section{Predicted KPIs and Safety Semantics}
\label{sec:dt_kpis}

A single evaluation returns a structured result containing the offered load,
predicted \code{power\_watts}, \code{delay\_us}, \code{throughput\_gbps}, and
\code{predicted\_loss}, together with two boolean flags that encode the twin's
two distinct concerns:

\begin{description}
  \item[\code{is\_safe} (QoS preserved).] True iff the predicted downlink delay
    is within the configured budget \emph{and} the predicted packet loss does
    not exceed its tolerance:
    \begin{equation}
      \texttt{is\_safe} \;=\;
      \bigl(\text{loss} \le L_{\max}\bigr) \;\wedge\;
      \bigl(\text{delay} \le D_{\max}\bigr),
      \label{eq:dt_is_safe}
    \end{equation}
    with the default budget $D_{\max} = 200~\mu$s and loss tolerance
    $L_{\max} = 5$ packets per interval. \dpdk{} is safe across the entire
    measured load range; \usr{} is safe only below its knee. The tolerance is
    deliberately non-zero: it is shared with the controller work of
    Chapter~\ref{chap:controllers}, where it additionally serves as the
    normalisation base of a graded QoS score. Its value is not cosmetic---as
    Section~\ref{sec:dt_thresholds} shows, it moves the derived QoS limit and
    thereby determines which constraint binds the decision threshold.
  \item[\code{is\_efficient} (locally cheaper).] True iff this action's
    predicted power is lower than the \emph{other} action's predicted power at
    the same load. It is a relative, per-load comparison---not an absolute
    threshold---and is the signal a controller uses to decide whether the
    energy-cheaper action is still the safe one.
\end{description}

The two flags are intentionally orthogonal. A \usr{} result with
\code{is\_safe}=True and \code{is\_efficient}=False denotes a load at which
\usr{} still meets QoS but already costs more than \dpdk{}: the regime in
which a controller should prefer \dpdk{} on energy grounds even though \usr{}
would not violate QoS.

Figure~\ref{fig:dt_surrogate} shows the surrogate's predicted responses over
the deployable load range. The contrast from Chapter~\ref{chap:upf_profiling}
is reproduced by the twin: \dpdk{} power and CPU are essentially flat, while
\usr{} rises with load and its delay and loss turn sharply upward past the
knee. The piecewise-constant ``staircase'' in the \usr{} curves is the
signature of the tree-based lite models; it is bounded to the measured
envelope and does not extrapolate, which is the desired conservative
behaviour at the edge of the data.

\begin{figure}[htbp]
  \centering
  \includegraphics[width=\linewidth]{\figpath{fig_dt_surrogate_curves.png}}
  \caption{Twin-predicted UPF responses versus offered load (lite surrogate).
    (a)~Power: \dpdk{} flat near 0.82\,W, \usr{} rising and crossing \dpdk{}
    at low load; (b)~CPU; (c)~downlink delay (log scale) with the 200\,$\mu$s
    QoS budget marked; (d)~predicted packet loss (symmetric-log scale). The
    \usr{} staircase reflects the piecewise-constant nature of the tree-based
    lite models.}
  \label{fig:dt_surrogate}
\end{figure}

% -- Section 6: Switching-cost model ------------------------------------
\section{Switching-Cost Model}
\label{sec:dt_switching}

Switching from one \upf{} to the other is not free: the new data plane must be
activated, which takes time and energy. Modelling this cost is essential, both
to penalise controllers that flip excessively and to represent the brief
window during which the new \upf{} is not yet serving traffic.

\subsection{Activation duration and the energy-spike rule}

The activation \emph{duration} is a property of the software stack, not of the
host: bringing up the DPDK poll-mode driver takes substantially longer than
starting the interrupt-driven \usr{} path. The twin uses measured durations of
24.0\,s for \dpdk{} and 3.3\,s for \usr{}.

Absolute switching \emph{energy} measured on one machine is not portable to
another, but the relationship
\begin{equation}
  E_{\mathrm{spike}}(\lambda) \;=\;
  \frac{P_{\mathrm{steady}}(\lambda)\,\cdot\,t_{\mathrm{act}}}{3600}\quad[\text{Wh}]
  \label{eq:dt_spike}
\end{equation}
holds within the source measurements and \emph{is} portable, because it
combines a target-machine quantity (the surrogate's predicted steady-state
power $P_{\mathrm{steady}}$) with a software-stack quantity (the activation
duration $t_{\mathrm{act}}$). The rule is validated against the source data:
for \dpdk{}, $34.81~\text{W} \times 24~\text{s} / 3600 \approx 0.232$\,Wh,
matching the measured spike; for \usr{}, $7.45~\text{W} \times 3.3~\text{s}
/ 3600 \approx 0.007$\,Wh.

\subsection{Within-step accounting}

Because a simulation step (15\,min at NetMob resolution) is far longer than
any activation, the twin must decide how a switch that occurs \emph{inside} a
step is accounted for. Three policies are supported
(Figure~\ref{fig:dt_switching}); the default is \emph{sub-step}:
\begin{description}
  \item[\code{sub\_step} (default).] The old \upf{} serves for
    $t_{\mathrm{act}}$ seconds and the new \upf{} for the remainder; the
    step's power is the time-weighted average, and the energy spike is paid
    once. KPIs are attributed to whichever \upf{} served the majority of the
    step, while safety is treated conservatively (a step is safe only if
    \emph{both} the old and new \upf{} are safe at that load).
  \item[\code{round\_down}.] The switch is treated as instantaneous within the
    step; only the spike is paid.
  \item[\code{round\_up}.] The new \upf{} is treated as unavailable for the
    entire step (worst case); the requested action does not take effect until
    the next step.
\end{description}

An optional \emph{prewarm} mode keeps \dpdk{} in standby so that switching to
it is effectively instantaneous, at the cost of a continuous standby power
while it is not the serving \upf{}. Prewarm is \emph{disabled} in the scenario
reported here, so \dpdk{} cold-starts and the activation cost above applies in
full; no standby draw is charged. This matches the deployment assumption of
Chapter~\ref{chap:controllers} and yields the more conservative claim, in that
no energy is attributed to a standby the deployment may not actually pay.

\begin{figure}[htbp]
  \centering
  \drawioplaceholder{%
    \textbf{Within-step switching accounting (\texttt{sub\_step}).} A single
    horizontal time bar representing one 15-minute step. Partition it into a
    short leading segment of width $t_{\mathrm{act}}$ (label ``old UPF
    serving'') and a long trailing segment (label ``new UPF serving''). Above
    the bar, draw the resulting step power as a weighted average
    $\bar P = f\,P_{\mathrm{old}} + (1-f)\,P_{\mathrm{new}}$ with
    $f = t_{\mathrm{act}}/t_{\mathrm{step}}$. Mark a one-off ``energy spike
    $E_{\mathrm{spike}}$'' at the boundary. Optionally show the two
    alternatives as smaller bars: \texttt{round\_down} (spike only, instant
    switch) and \texttt{round\_up} (entire step still on the old UPF, switch
    deferred). Keep the same colour convention as
    Figure~\ref{fig:dt_surrogate} (DPDK blue, USR red).}
  \caption{Within-step accounting of a switch under the \code{sub\_step}
    policy. The step power is the time-weighted average of the old and new
    \upf{} power; the activation energy spike is paid once.}
  \label{fig:dt_switching}
\end{figure}

% -- Section 7: Derived operating thresholds ----------------------------
\section{Operating Thresholds Derived from the Surrogate}
\label{sec:dt_thresholds}

A naive controller might switch \upf{}s at a hand-tuned load threshold. The
twin instead \emph{derives} its operating points from the surrogate itself, so
that they remain consistent with the measured physics and update automatically
if the models are retrained. Two quantities are obtained by sweeping the
stateless twin over load (Figure~\ref{fig:dt_threshold}):

\begin{description}
  \item[Energy break-even $\lambda_{\mathrm{be}}$.] The lowest load at which
    \usr{} power meets or exceeds \dpdk{} power. Below it, \usr{} is the
    energy-cheaper action; above it, \dpdk{} is.
  \item[QoS limit $\lambda_{\mathrm{qos}}$.] The highest load at which \usr{}
    is still \code{is\_safe}. This is the point the sweep \emph{can only find
    because the twin uses \code{usr\_full}}: it is where the surrogate's
    predicted loss or delay first crosses the budget.
\end{description}

\noindent
For the scenario studied here the sweep yields
$\lambda_{\mathrm{be}} = 91$\,Mbps and $\lambda_{\mathrm{qos}} = 149$\,Mbps. The
decision threshold a controller should use is the more conservative of the
two, minus a safety margin:
\begin{equation}
  \lambda_{\mathrm{dec}} \;=\;
  \min(\lambda_{\mathrm{be}}, \lambda_{\mathrm{qos}}) - m,
  \qquad m = 10~\text{Mbps},
  \label{eq:dt_decision}
\end{equation}
giving $\lambda_{\mathrm{dec}} = 81$\,Mbps here (energy-limited: the
break-even binds before the QoS limit). It is worth noting that this ordering
depends on the loss tolerance $L_{\max}$ of
Equation~\eqref{eq:dt_is_safe}. Under zero tolerance the QoS limit falls to
81\,Mbps and binds first, giving $\lambda_{\mathrm{dec}} = 71$\,Mbps and a
QoS-limited regime. The derivation machinery is unchanged either way---which
is the point of deriving rather than hardcoding---but the reported operating
point is sensitive to a budget that is a policy choice rather than a measured
quantity.

\subsection{Anti-flapping from forecast uncertainty}

A single decision threshold applied to a noisy forecast would cause rapid,
wasteful flipping near the boundary. The twin derives a hysteresis band
directly from the \emph{forecaster's} reported error, so that a single noisy
forecast cannot trigger a switch:
\begin{equation}
  \text{band} = 2 \cdot \mathrm{MAE}_{\mathrm{forecast}},\qquad
  \lambda_{\uparrow} = \lambda_{\mathrm{dec}},\qquad
  \lambda_{\downarrow} = \lambda_{\mathrm{dec}} - \text{band}.
  \label{eq:dt_band}
\end{equation}
The controller switches \usr{}$\to$\dpdk{} only when the predicted load rises
above $\lambda_{\uparrow}$, and back only when it falls below
$\lambda_{\downarrow}$.

This error-derived rule, however, has a failure mode that the present scenario
exposes and that is worth reporting explicitly. The forecaster's error is
reported in normalised units, so the band inherits the scale factor $\alpha$
of Equation~\eqref{eq:dt_alpha}. At $\alpha = 1$ the mean absolute error for
the operating cluster count is $\approx 110.7$\,Mbps, so
Equation~\eqref{eq:dt_band} would prescribe a band of $\approx 221$\,Mbps---
substantially \emph{wider} than the decision threshold itself. The lower
switch point is then clamped at $\lambda_{\downarrow} = 0$, the controller can
never return to \usr{} once it has left, and hysteresis silently degenerates
into a static-\dpdk{} policy. It is a benign-looking degeneration: the
controller still runs, still reports metrics, and simply never switches.

Two conclusions follow. First, an anti-flapping band derived from forecast
error is only meaningful when that error is small relative to the operating
point; when the forecaster is noisy at the scale of the decision threshold,
the rule as stated is unusable and must be capped rather than applied
literally. The twin therefore emits an explicit warning whenever the derived
band meets or exceeds the decision threshold. Second, for the results reported
here the band is instead fixed at
\begin{equation}
  \text{band} = 20~\text{Mbps},\qquad
  \lambda_{\uparrow} = 81~\text{Mbps},\qquad
  \lambda_{\downarrow} = 61~\text{Mbps},
  \label{eq:dt_band_used}
\end{equation}
which is narrow enough to preserve genuine two-threshold behaviour and matches
the band used for the classical baselines of
Chapter~\ref{chap:controllers}. Equation~\eqref{eq:dt_band} is retained as the
principled default for regimes in which the forecast error is small compared
with $\lambda_{\mathrm{dec}}$.

\begin{figure}[htbp]
  \centering
  \includegraphics[width=\linewidth]{\figpath{fig_dt_threshold_derivation.png}}
  \caption{Operating thresholds derived from the surrogate. \usr{} power
    (red) crosses \dpdk{} power (blue) at the energy break-even
    ($91$\,Mbps); the QoS limit ($149$\,Mbps) is where \usr{} first becomes
    unsafe. The decision threshold $\lambda_{\uparrow}=81$\,Mbps is the
    binding (energy break-even) limit minus a 10\,Mbps margin; the shaded band
    down to $\lambda_{\downarrow}=61$\,Mbps is the hysteresis region.
    The green region marks loads where \usr{} is both safe and cheaper than
    \dpdk{}.}
  \label{fig:dt_threshold}
\end{figure}

% -- Section 8: Evaluation methodology + demonstration ------------------
\section{Evaluation Methodology and Demonstration}
\label{sec:dt_eval}

\subsection{Rollout engine}

To evaluate any controller, the twin is driven over the full test set of
per-cluster load windows. For efficiency the rollout pre-evaluates the
stateless twin once for both actions over the entire flattened array of loads,
then performs the conceptually sequential per-cluster loop in pure Python:
each step reads the controller's requested action, looks up the precomputed
KPIs, and applies the switching-cost accounting of
Section~\ref{sec:dt_switching}. The result is one record per cluster-step,
suitable for aggregation. In the scenario studied here this is $K=10$ clusters
$\times$ $1009$ windows $= 10{,}090$ steps of 15\,min each.

\subsection{Metrics}

From the per-step records the following metrics are computed. Energy includes
both steady-state consumption ($\text{power} \times \Delta t$) and the
switching spikes of Equation~\eqref{eq:dt_spike}, so that flipping is
penalised:
\begin{description}
  \item[Total energy (Wh)] and \textbf{average power (W)} over all steps.
  \item[Energy saving (\%)] relative to the always-\dpdk{} baseline.
  \item[Unsafe USR rate] the fraction of \usr{}-serving steps that violated
    QoS---the key safety metric.
  \item[USR usage (\%)] the fraction of steps served by \usr{}.
  \item[Decision flip rate] the per-cluster fraction of steps at which the
    serving \upf{} changed---a proxy for churn.
\end{description}

\subsection{Controllers used in this chapter}

As stated in Section~\ref{sec:dt_intro}, this chapter exercises only
non-learning controllers, so that the figures characterise the \emph{twin}
rather than any particular learned policy:
\begin{description}
  \item[Static \dpdk{}] always selects \dpdk{} (the safe, high-power
    reference).
  \item[Static \usr{}] always selects \usr{} (maximally energy-greedy, and
    unsafe whenever load exceeds the knee).
  \item[Hysteresis] the rule of Equation~\eqref{eq:dt_band}, consuming the
    surrogate-derived thresholds and the forecast-derived band, with a
    one-step cooldown after each switch.
\end{description}
Learned controllers, an informed oracle upper bound, and the single-threshold
policy are evaluated in Chapter~\ref{chap:controllers}; they are deliberately
omitted here.

\subsection{How much room is there for energy-aware switching?}

The case for switching at all rests on the load distribution.
Figure~\ref{fig:dt_load_dist} shows that the offered load is heavily
right-skewed, but that under the calibration of
Equation~\eqref{eq:dt_alpha} only about 14\,\% of cluster-steps fall below the
decision threshold---that is, only on roughly one step in seven is \usr{} both
safe and cheaper than \dpdk{}. The remaining steps sit above the \usr{}
saturation knee, where \dpdk{} is the only viable action.

This bounds the achievable benefit tightly, and it is more honest to establish
that bound before presenting any controller. An informed oracle, given the
\emph{actual} load at each step and choosing the cheaper safe action, recovers
6.7\,\% of static \dpdk{}'s energy. No causal controller operating on this
traffic, this hardware, and this calibration can do better. The energy-aware
switching problem, in this regime, is a narrow one: the value of a controller
lies in capturing as much of that 6.7\,\% as possible without incurring QoS
violations, rather than in a large absolute saving.

The reason is physical rather than algorithmic. The profiled \usr{} data path
saturates at roughly $91$--$149$\,Mbps (Section~\ref{sec:dt_thresholds}),
whereas $\alpha = 1$ places cluster loads at $0.11$--$1.66$\,Gbps mean with
peaks to $5.8$\,Gbps---one to two orders of magnitude above the knee. A
software \upf{} of this capacity is simply undersized relative to the offered
traffic, and the twin's role here is precisely to make that mismatch
quantitative rather than to disguise it.

\begin{figure}[htbp]
  \centering
  \includegraphics[width=\linewidth]{\figpath{fig_dt_load_distribution.png}}
  \caption{Distribution of calibrated offered load across all clusters and
    time windows. Only about 14\,\% of cluster-steps lie below the derived
    decision threshold (dashed), i.e.\ in the regime where \usr{} is the safe
    and cheaper action; the long right tail extends well beyond the \usr{}
    saturation knee, where \dpdk{} is the only viable action.}
  \label{fig:dt_load_dist}
\end{figure}

\subsection{The twin under a hysteresis controller}

Figure~\ref{fig:dt_rollout} shows a representative window of one cluster's
rollout. The controller holds \usr{} while the predicted load stays within the
hysteresis band, switches to \dpdk{} when load rises above $\lambda_{\uparrow}$,
and returns to \usr{} only after load falls below $\lambda_{\downarrow}$; the
power trace (lower panel) accordingly sits near \dpdk{}'s 0.82\,W plateau
during high-load excursions and drops into \usr{}'s low-power regime
otherwise. The band visibly prevents single-sample flips around the threshold:
the load crosses $\lambda_{\uparrow}$ repeatedly without provoking a switch
until it also clears the band. Because prewarm is disabled, the cold-start
activation cost is now visible directly in the trace as short spikes above the
\dpdk{} plateau at switch instants.

\begin{figure}[htbp]
  \centering
  \includegraphics[width=\linewidth]{\figpath{fig_dt_hysteresis_rollout.png}}
  \caption{A representative window of the hysteresis controller driving the
    twin (one cluster). Upper: actual and predicted load with the switch
    points $\lambda_{\uparrow}$, $\lambda_{\downarrow}$; shaded spans mark
    \usr{}-served steps. Lower: the twin's predicted power, near the \dpdk{}
    floor when \dpdk{} serves and low when \usr{} serves.}
  \label{fig:dt_rollout}
\end{figure}

\subsection{Aggregate behaviour}

Table~\ref{tab:dt_metrics} and Figure~\ref{fig:dt_comparison} summarise the
controllers over all $10{,}090$ steps. The headline observation is that static
\usr{} is not merely unsafe here but actively \emph{more expensive}: it
consumes $2.5\times$ static \dpdk{}'s energy ($-154.5$\,\% ``saving'') while
violating QoS on 70.3\,\% of its steps. This is the direct consequence of
running a saturated data path---past its knee, \usr{} loses throughput,
inflates delay, and burns more CPU power than the accelerated path it was
meant to undercut. Any intuition that the software \upf{} is the cheap option
holds only below the knee.

The hysteresis rule recovers 4.8\,\% of static \dpdk{}'s energy while keeping
the unsafe rate at 3.0\,\%, by serving \usr{} on the 11.5\,\% of steps where
it is genuinely both safe and cheaper and ceding to \dpdk{} elsewhere. Against
the 6.7\,\% oracle bound established above, it captures roughly 72\,\% of the
attainable headroom using only a forecast and two derived thresholds. The
stateless threshold policy achieves a marginally higher saving (4.9\,\%) at a
noticeably worse unsafe rate (4.2\,\%) and twice the switching churn, which is
the trade the hysteresis band exists to make.

The absolute saving is modest, and it would be misleading to present it
otherwise: as established above, the regime admits at most 6.7\,\%. The
demonstration this chapter offers is therefore not that switching saves a
great deal of energy on this hardware, but that the twin \emph{quantifies the
trade-off correctly}---it identifies where the saturation knee lies, prices
the switching cost, bounds the attainable benefit, and exposes the
energy--safety frontier that learned controllers will later be trained to
navigate. A twin that reported a large saving here would be concealing the
capacity mismatch rather than measuring it.

\begin{table}[htbp]
  \centering
  \caption{Aggregate metrics over $10{,}090$ cluster-steps (non-learning
    controllers; learned controllers are deferred to
    Chapter~\ref{chap:controllers}). Energy includes switching spikes. The
    informed oracle sees the actual load at each step and is reported as an
    upper bound, not as a realisable controller.}
  \label{tab:dt_metrics}
  \setlength{\tabcolsep}{6pt}
  \begin{tabular}{lrrrrr}
    \toprule
    \textbf{Controller} & \textbf{Energy (Wh)} & \textbf{Avg power (W)}
      & \textbf{Saving (\%)} & \textbf{Unsafe USR (\%)} & \textbf{USR use (\%)} \\
    \midrule
    Static \dpdk{} & $2070.7$ & $0.821$ & $0.0$    & $0.0$  & $0.0$   \\
    Static \usr{}  & $5269.8$ & $2.089$ & $-154.5$ & $70.3$ & $100.0$ \\
    Threshold      & $1969.6$ & $0.780$ & $4.9$    & $4.2$  & $14.0$  \\
    Hysteresis     & $1971.0$ & $0.781$ & $4.8$    & $3.0$  & $11.5$  \\
    \midrule
    Oracle (bound) & $1931.0$ & $0.765$ & $6.7$    & $0.0$  & $13.7$  \\
    \bottomrule
  \end{tabular}
\end{table}

\begin{figure}[htbp]
  \centering
  \includegraphics[width=\linewidth]{\figpath{fig_dt_policy_comparison.png}}
  \caption{Aggregate comparison of the non-learning controllers.
    (a)~Total energy (with saving vs.\ static \dpdk{} annotated);
    (b)~unsafe \usr{} rate; (c)~\usr{} usage. Static \usr{} is both the most
    expensive and by far the least safe option in this regime, since it runs a
    saturated data path; hysteresis achieves a modest saving while keeping the
    unsafe rate low.}
  \label{fig:dt_comparison}
\end{figure}

% -- Section 9: Summary -------------------------------------------------
\section{Summary}
\label{sec:dt_summary}

This chapter constructed a measurement-grounded digital twin of the \upf{}
data plane from the profiling results of Chapter~\ref{chap:upf_profiling}. The
principal points are:

\begin{enumerate}
  \item \textbf{Scope.} The twin answers ``what would happen'' for a given
    (load, action) pair; it is neither a forecaster
    (Chapter~\ref{chap:forecaster}) nor a controller
    (Chapter~\ref{chap:controllers}), and is kept strictly controller-agnostic.
  \item \textbf{Variant choice.} \code{USR} is represented by \code{usr\_full},
    not \code{usr\_safe}, precisely because a safety-aware twin must be able to
    predict the saturated regime; \code{usr\_safe} would hide the failure mode
    and break threshold derivation.
  \item \textbf{Two modes.} A stateless physics mode and a stateful rollout
    mode share one surrogate core; the latter adds a portable switching-cost
    model based on a measured activation duration and an energy-spike scaling
    rule.
  \item \textbf{Derived, not hardcoded, thresholds.} The energy break-even and
    QoS limit are swept from the surrogate, and an anti-flapping band is
    derived from the forecaster's own error, keeping operating points
    consistent with the measured physics.
  \item \textbf{Demonstration.} Driven by a transparent hysteresis controller,
    the twin recovers 4.8\,\% of \dpdk{}'s energy---roughly 72\,\% of the
    6.7\,\% available to an informed oracle---while remaining safe on 97.0\,\%
    of \usr{} steps, quantitatively exposing the energy--safety trade-off that
    learned controllers will later be trained to navigate.
  \item \textbf{Bounded headroom.} Under the declared calibration
    $\alpha = 1$, the profiled \usr{} path saturates one to two orders of
    magnitude below the offered cluster loads, so only $\approx 14$\,\% of
    steps admit a \usr{} decision at all and static \usr{} costs $2.5\times$
    \dpdk{}. Establishing this bound is itself a result: it tells the
    controller chapter how much value a learned policy can possibly add, and
    it identifies \upf{} capacity---not control policy---as the dominant
    lever in this deployment.
\end{enumerate}

\end{document}
